% --- Archivo: sections/1_intro.tex ---
\section{Introducción}

El presente documento aborda el estudio de la dinámica ondulatoria cuántica, escalando desde la fenomenología unidimensional elemental hasta sistemas bidimensionales e interferencia de materia masiva. Inicialmente, se analiza el comportamiento de una partícula cuántica confinada en un pozo de potencial unidimensional con paredes impenetrables en $x=0$ y $x=L$, y su interacción con una barrera de potencial rectangular centrada en $x=x_c<L$. Clásicamente, si la energía cinética es inferior a la altura de la barrera, la transmisión es estrictamente nula. Sin embargo, su naturaleza cuántica y ondulatoria permite una probabilidad no nula de atravesarla, fenómeno conocido como efecto túnel. El objetivo de esta fase es resolver el sistema, visualizar la densidad de probabilidad interactuando con la barrera, y cuantificar las probabilidades de transmisión y reflexión de forma continua así como su dependencia con las características de la barrera.

Posteriormente, el alcance de la simulación se extiende al dominio bidimensional para modelar el paradigmático experimento de la doble rendija, visualizando la auto-interferencia espacial de un paquete de ondas. Finalmente, el proyecto da el salto a la frontera de la física moderna simulando la difracción de moléculas masivas de nanografeno (Hexabenzocoroneno, $C_{42}H_{18}$), evidenciando que la dualidad onda-corpúsculo persiste en escalas casi macroscópicas y justificando el uso de diferentes arquitecturas algorítmicas (Crank-Nicolson y métodos pseudo-espectrales) para su resolución.